% !TEX root = ../main.tex
%% Devoir surveillé n° 4 — chapitres 11 à 13 (valeurs vérifiées : verifications/devoirs.py)

\devoirDS{4}{Chapitres 11 à 13}{{La Station spatiale internationale : interaction gravitationnelle}/7,{La plongée sous-marine : fluide au repos}/7,{Le curling : forces et variation de vitesse}/6}

\begin{exoBac}{7}{La Station spatiale internationale}
La Station spatiale internationale (ISS) tourne autour de la Terre à une altitude d'environ \qty{400}{km}. À son bord, les astronautes
flottent en apesanteur.
\begin{donnees}
Constante de gravitation : $G=\qty{6,67e-11}{N.m^2.kg^{-2}}$. Terre : masse $M_\mathrm{T}=\qty{5,97e24}{kg}$, rayon $R_\mathrm{T}=\qty{6371}{km}$ ; intensité
de la pesanteur au sol : \qty{9,81}{N.kg^{-1}}. Masse de l'ISS : $m=\qty{4,20e5}{kg}$. Lune : masse $M_\mathrm{L}=\qty{7,35e22}{kg}$ ; on prendra
\qty{3,84e8}{m} pour la distance entre l'ISS et le centre de la Lune.
\end{donnees}
\begin{questions}
\item Donner l'expression de la valeur de la force exercée par la Terre sur l'ISS, puis la calculer.
\item Représenter cette force sur un schéma, avec l'échelle \qty{1}{cm} pour \qty{1,0e6}{N}.
\item Calculer la valeur du champ de gravitation terrestre à l'altitude de l'ISS. Comparer à sa valeur au sol.
\item Calculer la valeur de la force exercée par la Lune sur l'ISS et la comparer à celle exercée par la Terre.
\item \initiative Un élève affirme : « Dans l'ISS, les astronautes flottent parce qu'ils ne sont plus attirés par la Terre. »
Discuter cette affirmation.
\end{questions}
\end{exoBac}

\begin{exoBac}{7}{La plongée sous-marine}
Un plongeur descend en mer en emportant un ordinateur de plongée, qui mesure la pression absolue de l'eau.
\begin{docu}{1}{Les relevés de l'ordinateur de plongée}
{\centering\small\begin{tabular}{@{}lcccccc@{}}\toprule
Profondeur $h$ (m) & 0 & 5 & 10 & 15 & 20 & 25 \\
Pression $P$ (bar) & 1,01 & 1,52 & 2,02 & 2,52 & 3,03 & 3,53 \\\bottomrule
\end{tabular}\par}
\schema{ds4-pression}
\end{docu}
\begin{donnees}
$g=\qty{9,81}{N.kg^{-1}}$ ; $\qty{1}{bar}=\qty{1e5}{Pa}$. On considère l'air comme un gaz dont la température reste constante.
\end{donnees}
\begin{questions}
\item Interpréter, à l'échelle microscopique, la pression exercée par l'eau sur la paroi de l'ordinateur.
\item Montrer que les mesures sont en accord avec la loi fondamentale de la statique des fluides et en déduire la masse
volumique de l'eau de mer.
\item À l'aide de cette loi, calculer la différence de pression entre la surface et \qty{20}{m} de profondeur. Comparer aux relevés.
\item Un apnéiste remplit ses poumons de \qty{6,0}{L} d'air en surface, puis descend à \qty{20}{m}. Quel est alors le volume de l'air de ses
poumons ? Nommer la loi utilisée.
\item \initiative Un plongeur équipé d'une bouteille respire à \qty{20}{m} de l'air à la pression de l'eau. Pourquoi ne doit-il jamais
remonter en bloquant sa respiration ?
\end{questions}
\end{exoBac}

\begin{exoBac}{6}{Le curling}
Au curling, on lance une pierre de granit qui glisse sur la glace. On étudie son mouvement, dans le référentiel terrestre, après
qu'elle a été lâchée.
\begin{docu}{1}{Chronophotographie}
\schema{ds4-curling}
{\centering\small\begin{tabular}{@{}lccccccc@{}}\toprule
Point & M$_0$ & M$_1$ & M$_2$ & M$_3$ & M$_4$ & M$_5$ & M$_6$ \\
$x$ (m) & 0 & 1,40 & 2,70 & 3,90 & 5,00 & 6,00 & 6,90 \\\bottomrule
\end{tabular}\par}
\end{docu}
\begin{donnees}
Masse de la pierre : $m=\qty{19,0}{kg}$ ; $g=\qty{9,81}{N.kg^{-1}}$ ; durée entre deux positions : $\Delta t=\qty{0,50}{s}$.
\end{donnees}
\begin{questions}
\item Quelle est la nature du mouvement de la pierre ? Justifier.
\item Calculer les valeurs des vitesses $v_2$ et $v_3$ aux points M$_2$ et M$_3$.
\item Représenter $\vec v_2$, $\vec v_3$ et $\Delta\vec v=\vec v_3-\vec v_2$ (échelle : \qty{1}{cm} pour \qty{0,5}{m.s^{-1}}).
\item En déduire la direction, le sens et la valeur de la somme des forces exercées sur la pierre.
\item Faire l'inventaire des forces exercées sur la pierre. En déduire la valeur de la force de frottement exercée par la glace.
\end{questions}
\end{exoBac}
